\section{Описание}
Целью лабораторных работ по курсу «Информационный поиск» является разработка полноценной поисковой системы для тематического корпуса документов, извлечённых из открытых источников. В рамках проекта реализован полный цикл обработки данных: от сбора и предобработки до построения индексов, выполнения запросов и предоставления пользовательского интерфейса.

Система ориентирована на анализ вопросов и ответов с платформы \textbf{Stack Exchange}, в частности — с разделов \texttt{ai.stackexchange.com} и \texttt{datascience.stackexchange.com}. Эти источники выбраны благодаря их высокому качеству контента, технической направленности и актуальности для специалистов в области информационных технологий и анализа данных.

\subsection*{Этапы реализации}
Работа включает следующие ключевые компоненты:
\begin{itemize}
    \item \textbf{Сбор данных}: автоматизированный загрузчик извлекает архивы \texttt{.7z} с сайта Archive.org, распаковывает XML-файлы (\texttt{Posts.xml}) и выбирает только основные сообщения (вопросы) с достаточным объёмом текста (не менее 100 слов).
    \item \textbf{Хранение корпуса}: каждый документ сохраняется как отдельный текстовый файл в директории \texttt{corpus/}, а метаданные (URL, заголовок, хеш содержимого, время загрузки) — в JSON-файле с тем же идентификатором.
    \item \textbf{Токенизация}: текст приводится к нижнему регистру, разбивается на слова, удаляются все неалфавитно-цифровые символы. Токены короче двух символов отбрасываются.
    \item \textbf{Стемминг}: реализован алгоритм, вдохновлённый Porter Stemmer, для нормализации слов перед индексацией. Это позволяет объединять формы одного слова (например, \texttt{running} → \texttt{run}).
    \item \textbf{Построение индексов}: 
    \begin{itemize}
        \item Обратный индекс (\texttt{inverted\_index.bin}) — отображает термин в список ID документов, где он встречается.
        \item Прямой индекс (\texttt{forward\_index.bin}) — хранит информацию о каждом документе: ID, заголовок, URL.
    \end{itemize}
    Индексы сериализуются в бинарном виде для компактности и быстрого чтения.
    \item \textbf{Булев поиск}: реализована поддержка сложных запросов с операторами \texttt{\&\&} (AND), \texttt{||} (OR), \texttt{!} (NOT) и скобочной группировкой. Поиск выполняется за миллисекунды благодаря эффективному слиянию отсортированных списков документов.
    \item \textbf{Веб-интерфейс}: с помощью фреймворка Flask реализован веб-сервер, позволяющий вводить запросы через браузер, просматривать результаты и использовать постраничную навигацию.
\end{itemize}

\subsection*{Ключевые особенности}
\begin{itemize}
    \item Все компоненты написаны на языках \textbf{Python} (сбор данных, веб-интерфейс) и \textbf{C++} (токенизация, стемминг, индексация, поиск) без использования сторонних библиотек, кроме минимально необходимых (\texttt{std::vector}, \texttt{std::string}, \texttt{py7zr}, \texttt{Flask}).
    \item Система обеспечивает \textbf{идемпотентность}: повторная загрузка не дублирует данные, проверяется хеш содержимого.
    \item Используется \textbf{вежливое поведение} при скачивании: задержка между запросами, HEAD-запросы для проверки изменений.
    \item Поддерживается масштабируемость: система способна обрабатывать десятки тысяч документов с высокой скоростью поиска.
\end{itemize}

\subsection*{Статистика корпуса}
В результате работы сформирован корпус из \textbf{30\,000 документов} (по 15\,000 на источник), средний размер документа — около \textbf{500–2\,000 слов}. Документы прошли полную очистку: удалены HTML-теги, мусорные символы, сохранён только семантически значимый текст.

\pagebreak
